\section{The conference abstract, read against this report}
\label{appendix:abstract}

The companion conference paper's abstract was registered on 22 September
2026 and is fixed. This appendix takes it sentence by sentence and says where
this report substantiates each statement, with the precision the retained
evidence allows, so that no sentence of the abstract is read further than the
report supports. Statements about product features that this report neither
models nor measures are marked as such and located in the product repository
at commit \code{252216ba}, the snapshot this report transcribes. Two facts
postdate the abstract and appear only here: the reached-set attribution
defect of the measured builds (\cref{ob:attribution}) and the specification
entry's wrong date for the LLVM pair (\cref{sec:eval:rq5:llvm}).

\begin{center}
\small
\begin{longtable}{@{}L{4.3cm}L{2.5cm}L{7.3cm}@{}}
\caption{Each sentence of the conference abstract, where the report
substantiates it, and the precise statement.}
\label{tab:abstract-map}\\
\toprule
Abstract & Where & Precise statement and caveat \\
\midrule
\endfirsthead
\toprule
Abstract & Where & Precise statement and caveat \\
\midrule
\endhead
``An npm engine feeds a Docker base; another image needs that published base
and an npm SDK. If the final image fails, its release must remain pending
after its providers succeed.'' & \cref{sec:model} (\cref{def:window}),
\cref{thm:g2}; \cref{sec:eval:rq3}, \cref{tab:eval:rq3} & The
consumer-relative window keeps a contribution admissible until the consumer
releases it (\cref{thm:g2}). The three- and four-package registry topologies
of the companion campaigns exercise the npm-to-Docker chain: per topology, 18
designed failures and 12 repairs at unchanged planned versions, on one Linux
ARM64 Docker host. Which consumer waits for a publication is the declared
edge's matter: a Docker consumer built in an isolated registry-based builder
needs the published base, a consumer that reads local output does not
(\cref{sec:exec-durations}). \\
``We present a Git-based protocol that tags completed releases and resumes
unfinished work through the same release command.'' & \cref{sec:protocol};
\cref{thm:resumption}; \cref{sec:eval:rq2} & Recorded legs are annotated tags
and recovery is re-execution of the same command (\cref{thm:resumption}).
Observed: 20 of 20 recovery episodes converged within four runs in each
companion campaign, and 130 of 130 repeated-fault rows of the mechanism matrix
preserved plan and tags (\cref{tab:eval:mechanism}). The theorems concern the
model under the retry invariant; the implementation's departures are the
ledger of \cref{sec:obligations-current}. \\
``Folders define packages and shell commands define stages across languages,
single-package repositories, monorepos and polyrepos.'' &
\cref{sec:exec-polyrepo}; product documentation & The report models the
release computation and transcribes the three history modes of \spec{27}. The
configuration language (folders as packages, stages as commands) is the
product's, documented under \evpath{packages/docs/docs/configuration/} at
\code{252216ba}, and is not modelled here. No campaign of
\cref{sec:evaluation} evaluates a polyrepo or discovered-history release; those
semantics are transcribed, not evaluated, and the theorems of
\cref{sec:guarantees} are stated for a single history. \\
``Commits select releases, holds, prereleases, dependent updates and tests; an
agent skill supports diagnostics and automation.'' & \cref{sec:model},
\cref{sec:propagation}; \cref{sec:eval:rq3}; product
\evpath{specs/agent-guide/README.md} & Units, scopes, depth, channels and holds
are the model of \cref{sec:model,sec:propagation}; test selection by commit is
the CI fixture of \cref{sec:eval:rq3} (three repetitions, 81 predicates). The
agent guide (version 1.11.0-rc.3 at \code{252216ba}) is a product document for
coding agents; this report neither models it nor measures its effect, and no
campaign measures agent effectiveness. \\
``Shared Windows/Linux configuration, containerized tools and extensible
manifest readers let provisioned CI reuse its workflow as packages are
added.'' & product documentation only & Outside this report's model and
campaigns. Shared platform configuration, containerised tool execution and the
manifest scanner and writer are product features documented at
\code{252216ba} (\evpath{packages/docs/docs/configuration/},
\evpath{packages/docs/docs/cli/scanner.md}, \evpath{cli/writer.md}). No
campaign measures onboarding effort, and native Windows execution is
unmeasured everywhere (\cref{sec:eval:threats}). \\
``Terraform integration reconstructs temporary state from known resources
before dependent deployments.'' & \cref{sec:protocol}, the recorded-success
remark; product \evpath{packages/docs/docs/examples/terraform.md} & The
protocol section states what a Terraform-backed package's tag records and does
not: that a configuration version completed; it is not Terraform resource
state and does not prove that every live resource was imported. No campaign
of \cref{sec:evaluation} exercises the integration. \\
``Identical inputs produce identical, cacheable plans; the distributed profile
shares configuration and transports sources and results through Git.'' &
\cref{sec:determinism}; \cref{sec:complexity}; \cref{sec:exec-distributed};
\cref{sec:eval:rq5} & Determinism and idempotency are proved for the model
(\cref{sec:determinism}); the cache-reuse and optimisation claims of
\spec{13.11} and \spec{20.7} are restated with their qualifications in
\cref{sec:complexity}; placement invariance of the plan is the invariance
\cref{sec:exec-distributed} states for \spec{28}. Two releases through
Git-mailbox workers were measured, once per side (\cref{sec:eval:rq5}). The
measured builds carried the attribution defect of \cref{ob:attribution}, which
does not affect single-source units. \\
``Our evaluated single-history model builds on sagas and version control;
retry assumptions preserve versions and avoid repeating recorded releases.'' &
\cref{sec:intro}; \cref{thm:g3,thm:g4,thm:g5,thm:g6} & Version stability, no
double ledger delivery, no widening and convergence are theorems of the model
under the explicit retry invariant (\cref{sec:guarantees}); they say nothing
unconditional about external effects, and convergence counts progress-making
runs only. Sagas are credited in \cref{sec:intro}. \\
``All 120 controlled fault cases recovered after repair.'' &
\cref{sec:eval:rq2}, \cref{tab:eval:mechanism} & The 120 fault episodes of the
mechanism matrix (five graphs, two channels, four targets, three boundaries),
each executed twice under the fault, repaired and rerun: 130 of 130 repairs
left an empty plan (120 fault plus 10 foreign-identity cases) and 140 of 140
final no-ops left plan and tags unchanged. A designed deterministic matrix,
not a sample; the after-write fault is an adapter response, not a crash. \\
``Discarding consumer work at provider completion would instead lose releases
in seven out of ten cases (analytical comparison, not a tool benchmark).'' &
\cref{sec:eval:rq1}, \cref{tab:eval:counterfactual} & 84 of the 120 fault
episodes (70\%), 111 obligations, computed on the completion tags the binary
produced, with no modified binary run and no competitor; the ten no-fault
controls lose nothing and the ten foreign-identity cases are reported
separately. The finite projection finds the same loss in 260 of 772 states. \\
``Registry, CI, replay and scale experiments include google-cloud-node and
selected mid-release commit conflicts.'' & \cref{sec:eval:rq3};
\cref{sec:eval:rq2} (replay, orphan and push scenarios); \cref{sec:eval:rq4}
& Registry: the two topologies of the companion campaigns. CI: the
three-repetition fixture. Replay: five releases each of
\code{google-cloud-node} (36 of 42 exact) and \code{conventional-changelog}
(7 of 64) under reconstructed layouts. Mid-release conflicts: the E5 cells of
the 1.8.2 campaign, a clean and a conflicting concurrent push on selected
interleavings (12 and 14 checks), not a systematic sweep. Scale: the eight
timing cells of \cref{tab:eval:rq4} on a capped, contended host. \\
``Module suites include scanner/writer fixtures from real projects; release CI
requires the full suite regardless of change scope and checks binaries before
publication and installations afterwards.'' & \cref{sec:conformance}; the
product's release workflow; evidence
\evpath{experiments/distributed-20260922/evidence/product-gates/} &
\cref{sec:conformance} maps this report's results to the integration tests
behind them. The rest is the product's release policy as written in
\evpath{.github/workflows/release.yml} at \code{252216ba}: a ``Full suite''
job runs every package's tests on every release, the publication job depends
on it, and the workflow's header describes the binary build, the publication
and the install check as gates on one another. The release candidate's gate
logs of 22 September (the sharded suite 6 of 6 on \code{272772cb};
integration-only coverage 94.0\% against a 95\% gate, with a dispatch input to
lower it; the TinyGo gate skipped by dispatch input) are copied under
\evpath{product-gates/}. This report did not audit the workflow or the scanner
and writer fixtures; it reports the policy as read. \\
``Native Windows and distributed execution are outside these measurements.
Linux and LLVM results will be presented in the report.'' &
\cref{sec:eval:rq5}; \cref{sec:eval:threats} & Presented in
\cref{sec:eval:rq5}: kernel farm round 1 (every candidate rejected), round 2
(hexagon 95.89 and 92.22~s against 111, m68k\_nommu 115.61 and 115.48~s
against 128, after one rejected candidate; the single node 122.62 and
138.01~s) and the LLVM matrix (27:48.94 against 41:12.02, once per side; the
negative case 26:47.69 against 20:16). One repetition per side, asymmetric
hardware, no general speedup claimed. Native Windows remains unmeasured. \\
\bottomrule
\end{longtable}
\end{center}
